\documentclass[10pt,a4paper]{article}
\usepackage{fontspec,amsmath,amssymb,graphicx}
\usepackage[margin=18mm]{geometry}
\setmainfont{Times New Roman}
\setlength{\parindent}{0pt}
\setlength{\parskip}{5pt}
\allowdisplaybreaks
\sloppy
\emergencystretch=2em
\begin{document}
\section*{BT\_3\_6\_3\_07 — Lời giải bổ sung}
\subsection*{Kiến thức và trình tự giải}
\textbf{Bước 1 — Chọn Kirchhoff (3D) hoặc Poisson (2D).} Với li độ đầu bằng không, vận tốc $\psi$, tốc độ $c$, đặt $\rho=ct$. Trong ba chiều,


\[
u(P,t)=\frac t{4\pi}\int_{S^2}\psi(P+\rho\omega)d\omega.
\]


Trong hai chiều,


\[
u(P,t)=\frac1{2\pi c}\int_{|Q-P|<ct}\frac{\psi(Q)}{\sqrt{c^2t^2-|Q-P|^2}}dQ.
\]


Với li độ $f$ khác không, thêm đạo hàm theo $t$ của chính biểu thức tương ứng khi thay $\psi$ bằng $f$. Với nguồn và dữ liệu đầu không, tích phân theo thời gian của nghiệm có vận tốc là nguồn tại từng thời điểm (Duhamel).

\textbf{Bước 2 — Tận dụng hình học và đối xứng.} Chỉ thị miền $D$ biến Kirchhoff thành $t$ nhân tỷ lệ diện tích mặt cầu nằm trong $D$. Chọn trục qua tâm miền và $P$, đặt $\mu=\cos\theta$, $d\omega=d\mu\,d\phi$. Nếu điều kiện không phụ thuộc $\phi$, tỷ lệ là $|E|/2$ với $E\subset[-1,1]$. Mỗi nửa không gian cho một bất đẳng thức tuyến tính theo $\mu$; quả cầu cho một ngưỡng chỏm. Miền hợp/rút bỏ cho tổng/hiệu nếu các phần không chồng diện tích.

\textbf{Bước 3 — Viết chính xác các khoảng hợp lệ.} Giải bất đẳng thức gốc, sau đó mới bình phương khi có căn. Chia $[-1,1]$ ở nghiệm phương trình biên, kiểm tra dấu ở trung điểm từng khoảng, loại nghiệm ngoại lai. Nếu còn góc phương vị, dùng độ dài tập $\cos\phi<v$ rồi còn một tích phân theo $\mu$. Với Poisson, tích phân theo một tọa độ cho hiệu arcsin; cắt các cận vào đĩa trước khi tính.

\textbf{Bước 4 — Tính nguồn bằng nguyên hàm thời gian.} Với quả cầu bán kính $R$ cách $P$ khoảng $r>0$, đặt


\[
K_R(r,t)=t\left[\frac{R^2-(r-ct)^2}{4rct}\right]_0^1,\quad J_R(r,t)=\int_0^tK_R(r,s)ds.
\]


Tách ba giai đoạn: cả mặt cầu nằm trong miền, giao một chỏm, không giao. Nguồn bật trên $[0,T]$ cho $J_R(r,t)-J_R(r,(t-T)_+)$. Trường hợp $r=0$ tính trực tiếp, tránh chia cho $r$.

\textbf{Bước 5 — Kiểm tra chuẩn hóa, mặt sóng và cực đại.} Nếu $\psi=1$ trên toàn không gian thì $u=t$, dùng để kiểm tra hệ số. Với dữ liệu bước nhảy, chỉ kết luận cổ điển ngoài các mặt sóng. Cực đại phải xét cả điểm chuyển nhánh và giá trị một phía; nêu rõ supremum nếu chỉ đạt theo giới hạn. Nếu dùng tích phân số, ghi tham số và cách kiểm tra hội tụ; đồ thị không tự chứng minh một cực đại toàn cục.
\subsection*{Lời giải áp dụng cho bài}
\textbf{Bước 1. Tích phân Poisson theo tung độ trước.} Với $c=2$, tại hoành độ $X$, đĩa ảnh hưởng có tung độ $P_y\pm h(X)$. Một dải $l<Y<v$ góp nguyên hàm $\arcsin((Y-P_y)/h)$, sau khi cắt hai cận vào $[-1,1]$. Các miền trong đề có cận tung độ tường minh theo $|X|$, nên chỉ còn một tích phân theo $X$.

\textbf{Bước 2. Hai đối xứng cần dùng.} Câu (a): trước $2t<98/\sqrt5$, cả đĩa quanh $(0,100)$ nằm trong miền $D_1$, vận tốc bằng $2$, nên $u=2t$. Câu (b): $\psi(Y,X)=-\psi(X,Y)$ và điểm quan sát $(100,100)$ bất biến khi đổi tọa độ. Đĩa và nhân Poisson cũng bất biến, nên tích phân đổi thành đối của chính nó: $u(100,100,t)=0$ với mọi $t$. Công thức tích phân dưới đây áp dụng cả cho các điểm khác khi vẽ đồ thị, nhưng tại điểm của câu (b) kết quả rút gọn là $0$.

\textbf{Bước 3. Tính các dải và dựng đồ thị.}

Poisson hai chiều: $u(P,t)=(2\pi c)^{-1}\int_{|Q-P|<ct}\psi(Q)/\sqrt{c^2t^2-|Q-P|^2}\,dQ$. Tích phân theo tung độ được thực hiện chính xác. Đặt $h(X)=\sqrt{c^2t^2-(X-P_x)^2}$,
\[B_X(l,v)=\arcsin\left[\frac{v-P_y}{h(X)}\right]_{-1}^1-\arcsin\left[\frac{l-P_y}{h(X)}\right]_{-1}^1.\]
(a) Với $P=(0,100),c=2$,
\[u=\frac1{4\pi}\int_{-2t}^{2t}\{2B_X(2+2|X|,\infty)-B_X(-\infty,\min(2-2|X|,0))+B_X(\max(2-2|X|,0),2+2|X|)\}dX.\]
(b) Với $P=(100,100),c=2$, tích phân $X\in[100-2t,100+2t]$: nếu $|X|>1$, tích phân tung độ là $B_X(-1,1)$; nếu $|X|<1$, là $-B_X(-\infty,-1)-B_X(1,\infty)$. Nhân $1/(4\pi)$. Đây là tích phân một chiều của hàm sơ cấp; không có kỳ dị căn ở tung độ nữa. Ký hiệu $D_1$ thứ hai trong đề là $D_2$.
(c) Đồ thị kèm theo được tính từ tích phân Poisson; mã cho phép đổi điểm/thời gian. Dữ liệu chỉ thị nên hiểu nghiệm yếu và sự gãy ở biên vùng ảnh hưởng.

\textbf{Bước 4. Đối chiếu và kết luận.} Với vận tốc hằng $1$, tích phân hai chiều của nhân bằng $2\pi ct$, cho $u=t$, kiểm tra hệ số $1/(4\pi)$ khi $c=2$. Đồ thị bề mặt dùng tích phân số, phải đọc tham số lưới trong mã; đường $x=y$ của câu (b) luôn bằng $0$ theo đối xứng, là một kiểm tra độc lập cho đồ thị.

\textbf{Chi tiết loại tích phân tung độ và kiểm tra hệ số.} Tại $X$ cố định, đặt $h^2=c^2t^2-(X-P_x)^2$ và $z=(Y-P_y)/h$. Khi $h>0$,
\[\int\frac{dY}{\sqrt{h^2-(Y-P_y)^2}}
=\int\frac{dz}{\sqrt{1-z^2}}=\arcsin z.\]
Chỉ lấy giao của dải với $P_y-h<Y<P_y+h$, tương đương cắt $z$ vào $[-1,1]$, cho $B_X(l,v)$. Với vận tốc hằng $1$, dùng tọa độ cực quanh $P$:
\[\int_{|Q-P|<ct}\frac{dQ}{\sqrt{c^2t^2-|Q-P|^2}}
=2\pi\int_0^{ct}\frac{r\,dr}{\sqrt{c^2t^2-r^2}}
=2\pi ct.\]
Nhân $1/(2\pi c)$ cho $u=t$. Ở (a), khoảng cách từ $(0,100)$ đến hai đường $Y=2\pm2X$ là $|100-2|/\sqrt{1+4}=98/\sqrt5$, nên trước $2t<98/\sqrt5$ cả đĩa có vận tốc $2$.
\begin{center}\includegraphics[width=.97\linewidth,height=11cm,keepaspectratio]{Do_thi/BT_3_6_3_07.png}\end{center}
\textit{Đồ thị minh họa; với nghiệm số, xem phương pháp và tham số trong mã Python kèm theo.}\end{document}